%  LaTeX support: latex@mdpi.com 
%  For support, please attach all files needed for compiling as well as the log file, and specify your operating system, LaTeX version, and LaTeX editor.

%=================================================================
\documentclass[constrmater,article,accept,pdftex,moreauthors]{Definitions/mdpi} 

% MDPI internal commands
\firstpage{1} 
\makeatletter 
\setcounter{page}{\@firstpage} 
\makeatother
\pubvolume{1}
\issuenum{1}
\articlenumber{0}
\pubyear{2023}
\copyrightyear{2023}
%\externaleditor{\hl{Academic Editor:}Firstname Lastname}
\datereceived{19 January 2023} 
\daterevised{5 March 2023}
\dateaccepted{} 
\datepublished{} 
\hreflink{https://doi.org/} % If needed use \linebreak
%\doinum{}

%=================================================================
% Add packages and commands here. The following packages are loaded in our class file: fontenc, inputenc, calc, indentfirst, fancyhdr, graphicx, epstopdf, lastpage, ifthen, lineno, float, amsmath, setspace, enumitem, mathpazo, booktabs, titlesec, etoolbox, tabto, xcolor, soul, multirow, microtype, tikz, totcount, changepage, attrib, upgreek, cleveref, amsthm, hyphenat, natbib, hyperref, footmisc, url, geometry, newfloat, caption

\graphicspath{{figures/}} % path to folder with figures
\usepackage{subcaption}
\usepackage{listings}
\usepackage{xcolor}
\usepackage{gensymb} % for \textdegree
\usepackage{tabularx}
\usepackage{fixltx2e}
\newcommand{\highlight}[1]{\colorbox{yellow}{#1}} 
%=================================================================
% Full title of the paper (Capitalized)
\Title{Geotechnical Properties of Soil Stabilized with Blended Binders for Sustainable Road Base Applications}

% MDPI internal command: Title for citation in the left column
\TitleCitation{Geotechnical Properties of Soil Stabilized with Blended Binders for Sustainable Road Base Applications}

% Author Orchid ID: enter ID or remove command
\newcommand{\orcidauthorA}{0000-0002-0577-9936} % Per Add \orcidA{} behind the author's name
\newcommand{\orcidauthorB}{0000-0002-5759-1089} % Polina Add \orcidB{} behind the author's name

% Authors, for the paper (add full first names)
\Author{Per Lindh %MDPI: Please carefully check the accuracy of names and affiliations. <-- Authors' reply: yes, correct.
 $^{1,2,}$ %MDPI: The Correspondence author is different from our system. Please check which one is correct. <-- Authors' reply: corrected (Polina Lemenkova).
\orcidA{} and Polina Lemenkova $^{3,}$* %MDPI: We add this symbol here to keep the consistent with correspondence email and our system. Please confirm. <-- Authors' reply: yes, correct.
\orcidB{}}

%\longauthorlist{yes}

% MDPI internal command: Authors, for metadata in PDF
\AuthorNames{Per Lindh and Polina Lemenkova}

% MDPI internal command: Authors, for citation in the left column
\AuthorCitation{Lindh, P.; Lemenkova, P.}
% If this is a Chicago style journal: Lastname, Firstname, Firstname Lastname, and Firstname Lastname.

% Affiliations / Addresses (Add [1] after \address if there is only one affiliation.)
\address{%
$^{1}$ \quad Department of Investments, Technology and Environment, Swedish Transport Administration, Neptunigatan~52, Box 366, SE-201-23 %MDPI: The postcode of the same Country should be keep the consistent format. Please revise the affs 1 and 2. <-- Authors' reply: corrected to 'SE-201-23'.
 Malmö, Sweden \\
$^{2}$ \quad Division of Building Materials, Department of Building and Environmental Technology, Lunds Tekniska Högskola LTH (Faculty of Engineering), Lund University, Box 118, SE-221-00 Lund, Sweden \\
$^{3}$ \quad Laboratory of Image Synthesis and Analysis (LISA), École Polytechnique de Bruxelles (Brussels Faculty of Engineering), Université Libre de Bruxelles (ULB), %MDPI: We replace dot to comma. Please confirm. <-- Authors' reply: yes, we confirm.
 Building L, Campus du Solbosch, ULB--LISA CP165/57, Avenue Franklin D. Roosevelt 50, 1050 Brussels, Belgium
}

% Contact information of the corresponding author
\corres{Correspondence: %MDPI: Please confirm if this email should be retained. <-- Authors' reply: no, this email is deleted.
 or per.lindh@byggtek.lth.se (P.L.); polina.lemenkova@ulb.be (P.L.); Tel.: +46-0771-921-921 (P.L.); +32-471-86-04-59 (P.L.)}

% Current address and/or shared authorship
%\firstnote{Current address: Affiliation 3.} 

% Abstract (Do not insert blank lines, i.e. \\) 
\abstract{This study aimed at evaluating the effect of blended binders on the stabilization of clayey soils intended for use as road and pavement materials in selected regions of Sweden. The stabilization potential of blended binders containing five stabilizers (cement, bio fly ash, energy fly ash, slag and lime) was investigated using laboratory tests and statistical analysis. Soil samples were compacted using Swedish Standards on UCS. The specimens were stabilized with blended mixtures containing various ratios of five binders. The effects of changed ratio of binders on soil strength was analyzed using velocities of seismic P-waves penetrating the tested soil samples on the day 14 of the experiment. The difference in the soil surface response indicated variations in strength in the evaluated specimens. We tested combination of blended binders to improve the stabilization of clayey soil. The mix of slag/lime or slag/cement accelerated soil hardening process and gave durable soil product. We noted that pure lime (burnt or quenched) is best suited for the fine-grained soils containing clay minerals. Slag used in this study had a very finely ground structure and had hydraulic properties (hardens under water) without activation. Therefore, slag has a too slow curing process for it to be practical to use in real projects on stabilization of roads. The best performance on soil stabilization was demonstrated by blended binders consisted of lime / fly ash / cement which considerably improved the geotechnical properties and workability of soil and increased its strength. We conclude that bearing capacities of soil intended for road construction can be significantly improved by stabilization using mixed binders, compared to pure binders (cement).}

% Keywords
\keyword{soil strength; binders; road pavements; soil strength; recycling technologies; materials science; engineering; civil engineering; construction; compressive strength; soil stabilization}
%MDPI: We changed the first letter of all words to lowercase. Please check and confirm. <-- Authors' reply: yes, agree and confirm.

% The fields PACS, MSC, and JEL may be left empty or commented out if not applicable, from the layout guide, only keywords can be put here, please move  PACS, MSC, JEL parts to other parts. 
\PACS{81.40.Cd; 81.40.Ef; 62.20.Qp; 83.50.Xa; 45.70.Mg; 92.40.Lg; 81.40.Lm; 62.20.M–}
\MSC{76A99; 74E30; 74Fxx}%MDPI: Please check if these MSC code is correct. <-- Authors' reply: corrected to 76A99, 74E30 and 74F20.
\JEL{Q00; Q01; Q24; Q55; Q56}

%%%%%%%%%%%%%%%%%%%%%%%%%%%%%%%%%%%%%%%%%%
\begin{document}

\section{Introduction}

Selecting binders for soil stabilization is a widely recognized problem in civil engineering where one tries to find a proper type and percentage of stabilizing agents to improve soil properties~\cite{Consoli2019,daRocha,Freitag}. For~instance, using binders to improve the properties of soil is necessary for road construction. As~a result of stabilization by additives, such as cement, lime, fly ash and slag, soil obtains improved workability and better engineering characteristics and geotechnical properties. These are crucial for road construction due to the increased compressive strength, flexion strength, porosity, density and dynamic modulus~\cite{Carasca}. Many existing methods, such as~\cite{Mypati,Sargent}, deal with various binders and evaluate their effects on stabilization of various types of soil. Specifically, for~Sweden with its northern environmental and climate conditions stabilisation of expansive soil is crucial due to the strong effects on urban infrastructure \citep{Amorim,Gustavsson}. For~instance, stabilized soil should be more resistant to freeze--thaw cycles and fatigue cracking, which is important for transport network safety in cold~regions. 

Previous studies have indicated that the addition of various stabilizing agents improved the workability of soils \citep{Etim,Fasihnikoutalab,Kogbara,Lindh2022b,Mavroulidou}. For~instance, adding cement improves the strength and bearing capacities of soils~\cite{Ekpo}, while adding lime contributes to gain in strength in tested soil specimens~\cite{Frempong}. This suggests the possibility that the experimental addition of various stabilising agents can even more improve soil characteristics. Contrary to these methods, it was proposed recently to use mixtures of diverse combinations of blended binders. Thus, in~existing studies \citep{Abbey,AlHdabi,ma15217798,Jafer,Lietal2021,Wangetal2015} experimented with the addition of blended binders to soil which caused a significant gain in~strength.

Much has been written on the improved properties of soil after adding various binders as stabilizers. For example,~\cite{Consoli} reports that the addition of cement to the soft soils fastens acquiring strength results;~\cite{MESKINI2021129161} mentioned stabilizing effect of fly ash and lime additives on strength and durability of soil through neutralizing acidity;~\cite{Mirkovic} tested wearing course asphalt mixtures made with various percentages of fly ash and reported satisfactory volumetric composition which can be achieved by adding fly ash into soil mixture, improved bulk density and filled voids of the composed mineral and asphalt mixture;~\cite{app12189134} improved the mechanical behavior of an asphalt mixture using the alternative aggregate boiler coke ash, an~element that originates in nickel processing. Further;~\cite{Lima} evaluated the feasibility of using mud and fly ash residues to produce asphalt~mixtures. 

{Long-term practical development of soil stabilization in Sweden resulted in improved technical equipment used for field and laboratory soil investigations, and~new design methods} \cite{Bjerrum}. {Various binders of were tested and applied for a more effective soil stabilization. For~instance, }previous studies report~\cite{Higgins,electronics11223726} that the Ground Granulated Blast Furnace Slag (GGBFS) can be effectively used as a replacement of Portland cement in civil engineering and construction engineering works due to its high geotechnical quality and environmental sustainability. Further, the~long-term trend of the improvement of properties of the subgrade soils stabilized by fly ash was analyzed for several decades and reported in recent studies~\cite{Parsons}. Soil chemically stabilized with blends of cement and fly ash demonstrated improved characteristics in the Unconfined Compressive Strength (UCS) and split tensile strength tests, as~demonstrated in previous studies that also mentioned the increased shrinkage limit and crack reduction in soil samples~\cite{Olgun}.

However, while the results presented in existing studies on blended binders and their applications in various regions with diverse soil types~\cite{Jie,Holz,Jamnongpipatkul}, are consistent with those obtained by other works on effective methods of soil stabilization~\cite{Moh}, it is still unclear how these methods behave when used in specific conditions of Sweden with low winter temperatures and high technical requirements for soil and road specifications. {For instance, maintenance of roads during non-favourable conditions of winter-spring period in Sweden is based on using soil resisting to intense rainfalls and freeze--thaw cycles. Therefore, selecting effective binder combinations is especially important for stabilization of soil collected in southern Sweden.} As a response to such needs, the~goal of this study is to test how different components in binders in binder blends interact during {clayey }soil stabilization and increase the stability and safety of roads depending on proportions and ratio with regard to water content{ in clayey soil, which is essential for the climate setting of southern Sweden}. The~target improvements of soil specimens include gain in strength, resistivity and durability of the stabilized soil in roads, as~basic requirements for transport~safety.

{For instance, while the existing practices of soil stabilization focused on the design of }road constructions and transport engineering \citep{Celauro,Glossop,Markwick,Nelson,Olsenetal2022,LindhLemenkovaTTJ,Sear}, our study contributes to these investigations with a special aim at evaluating the effects of blended binders. {The use of blended binders }such as slag, cement, lime, energy fly ash and bio fly ash {enables to increase the quality of soil and thus to improve the road drainage measures in Swedish climate setting. Because~the combination of blended mixtures has better effects on soil properties, such experiments ensure more effective and stable road  construction, operation and maintenance. In~this way, effective }soil stabilization {increases} safety and durability of roads. {Further}, in~this paper we concentrate on presenting the mixed binders and evaluate their effects on the increase of soil properties intended as a sustainable road base. {Suggested improvements concerning using blended binders instead of single binders for soil stabilization included the statistical computing of the optimized ratio of the mixed binders.} Mixed binder {lead} to a better and more durable{ soil stabilized} for road pavement. On~the other hand, {we showed that }blended binders have different properties in terms of the reaction period with soil and curing time, which requires a special evaluation through a series of tests. In~this setting {we demonstrated that }hardening of soil {by} an optimal combination of blended binders {is} an important problem in civil engineering and construction~industry. 

{This aim of the present paper is to present current practice in stabilization of clayey soil used for building road surface in Sweden by blended binders. }To illustrate the practical importance of the blended binders for Swedish soils, we apply it to real tested clayey soil types collected from the sample area in southern Sweden. Clay is a fine-grained weak soil that tends to change in volume due to the fluctuations in water content under diverse weather conditions~\cite{Sudhakar}. Therefore, such type of soil should be stabilized prior to road construction to increase the strength and stability of road and transportation safety. There are many binder types used for stabilization of weak soil and reported for road construction cases \citep{Antunes,Park}. Portland cement and lime are the traditional binders in stabilization contexts where the lime needs clay mineral or air ($CO_2$) for a hardening process to take place \citep{Bell,Hilt,LindhLemenkovaCEE,Taylor}. 

{Previous} paper on stabilization of clays{ }exploited the importance and diverse effects of binders on clays~\cite{Okagbue,Estabragh,Xu}. {To continue this research}, we {extended the experiments on combinations of binders} for clays collected in southern Sweden and detect the effects from the percentage of various binders on hardening of clayey soil while trying to statistically balance the ratio of binders in a mixture. {Experimental trials with different binders used for soil stabilization in Swedish climate aimed at improving the hydraulic performance of soil such as long-term infiltration capacity and porosity.}

{Based on an experimental combination of various proportions of binders used for soil stabilization, the~study sought to identify the following aspects of road construction: (1) problems experienced concerning soil strength used for road construction, focusing on the current Swedish climate; (2) effective soil stabilization enables to avoid possible issues regarding climate impacts such as flash floods and heavy showers; and (3) adaptation measures regarding road stability taking environmental setting of Sweden into account, and~improved road functioning by increased strength parameters of clayey soil as a construction material and increased bearing capacity in roads.}

Aimed at {the }developing {of the} unified{ }framework of clayey soil stabilization by blended binders, our study continues the research on improvement its geotechnical properties which has been observed in multiple previous publications \citep{Sorengard,AlSwaidani,Feng,LindhLemenkovaECES,Maheepala,Lindh2022NL,Nordmark}. These and other works result from the need of the improved soil properties with the addition of binders in transport infrastructure, which can be achieved by the treatment of soils with stabilising agents. Positive results of increasing gain in strength upon stabilization with the addition of various binders are demonstrated in these papers. The~goal of this survey is to provide a series of tests aimed at improving the geotechnical properties of soil. The~objective is to ensure high engineering characteristics and safety of roads and sustainability of pavements{ in local settings of southern Sweden. Thus, the~overview of the location and condition of soil used for construction of road facilities was carried out to better select binders.} 

{Binders that work best for clayey soils of southern Sweden ensure stability and strength of roads during future exploitation. Suggestions concerning selecting binder blends, research designed based on simplex lattice experiments and seismic measurements used to evaluate P-wave velocity indicating soil strength aimed at the quality of the stabilized soil for improved road bearing capacity. }The economic effects from the stabilized soil has direct economic advantages in road construction industry such as higher material stability, decreased erosion, reduced road maintenance and construction costs, ensured safety in transport network to mention a few of them. All these factors together significantly contribute to the societal development and a well-being of the~society.

%%%%%%%%%%%%%%%%%%%%%%%%%%%%%%%%%%%%%%%%%%
\section{Materials and~Methods}

\subsection{Implementation of the~Experiment}

The following experiments were implemented to evaluate the performance of clayey soil after stabilization with five various binders. We compared the effects from cement, bio fly ash, energy fly ash, slag and lime on the gain in strength of soil specimens evaluated using P-wave velocities. In~the first step, soil specimens were fabricated using different binder recipes. To~ensure the significance level of the results, the~recipes are based on the previous experiments in combination with the statistical trial assessment. The~statistical planning minimizes the number of the experiments that need to be performed in order to reach the desired level of significance. For~five different binders, 100 test soil specimens were manufactured. These specimens were stored for 90 days to achieve a final strength even for the slow-curing components. During~the storage period, the~speed of the compression wave and shear wave have been measured in the material samples at different times to quantitatively measure the strength of soil~specimens. 

To investigate how five different binders work individually and in combination with each other, we used different approaches. A~common method is to keep all the parameters constant except for the target one that is varied. In~our case, the~target parameter was the amount of each tested binder: cement, bio fly ash, energy fly ash, slag and lime, respectively. This is repeated with the next parameter iteratively. Although~this method is conventionally used, it includes certain drawbacks as follows. First, the~interaction effects and the second-order effects between the different components are missed in the final evaluation of the soil strength. These interactions can be both positive and negative, i.e.,~missed interaction and interpretation of the results regarding safety issues. Therefore, the~methodology of our study considered more effective approaches, such as simplex-lattice design, while evaluating the results (Figure \ref{fig1}).

\begin{figure}[H]
%\centering
	\includegraphics[width=11.3 cm]{Fig_01.jpg}
\caption{Experimental design scheme. \textbf{Left}: 3$\times$2 factor experiment for evaluation of the effects of binders on soil strength. \textbf{Right}: Schematic image of a ``mixture design'' with three factors; A, B and C. The~experimental setup is called simplex-lattice (3$\times$2). \label{fig1}}
\end{figure}  
% please change Left and Right to a and b in the figure caption, and add a and b to the sub-figure.   The same as other figurs. 

As an alternative to the above, there is a response surface method. The~approach of the response surface method is based on the assumption that all the variables that affect the soil strength are varied according to a special scheme. In~a factorial trial, the~amount of binder is varied either according to the low and high admixture, or~triple admixture consisting of low, medium and high values depending on whether the second order effects are to be taken into account. Due to such an algorithm, this method is very effective for testing large amounts of soil (e.g., hundreds of tons of the materials), because~even few attempts provide enough data for decision making. Another advantage of the response surface method is that it enables to evaluate different binder levels. The~disadvantage is that in the specimens with high levels of several binders, the~content becomes very high which may affect the results. A~simple illustration with only two binders is indicated in Figure~\ref{fig1}.

\subsection{General~Workflow}

The design experiment of this study followed the existing guidance~\cite{Montgomery1997} and includes the following points in problem formulation and solving: 

\begin{itemize}
	\item	 Identifying the factors affecting soil strength;
	\item Defining the amounts and the ratio of binders optimal for effective stabilization of clayey soil; 
	\item Identifying the environmental limits affecting the geotechnical properties of soil;
	\item Selecting the response variables: cement, fly ash and lime in various ratios;
	\item Experimental modelling and statistical approaches for the data analysis;
	\item Performing the experiments using the UCS and P-wave seismic tests;
	\item Quality control aimed at identifying the advantages and drawbacks of the methods.
\end{itemize}

Parts of the above steps in the methodology take place iteratively, while others were performed in parallel. We considered the statistical approaches in data analysis for the experimental testing aimed at evaluating the optimal binder amount that should be added for the best results in soil stabilization \citep{Box}. To~increase the efficiency of the experiment, statistical data processing was applied and the outputs were visualized as ternary diagrams and Pareto charts. Moreover, we applied the additional points while performing the statistical data processing~\cite{Montgomery1997}: 

\begin{enumerate}
	\item	Considering the non-statistical knowledge to the geotechnical problems related to road construction; 
	\item Keeping the design and analysis as simple and clear as possible; 
	\item Identifying and recognition of the difference between the empirical practical and theoretical statistical significance; 
	\item Performing the iterative series of the experiments on soil specimens. 
\end{enumerate}	

	The methodology consists in the optimization of the process of soil stabilization the the main goal to evaluate the soil strength as a crucial criterion of safety for sustainable road~construction. 

\subsection{Manufacturing Soil~Specimens}
	The specimens were manufactured in a special {compaction} equipment {of compaction apparatus} that is designed to pack the stabilized fine-grained moraine~\cite{Lindh2004}, {Figure} \ref{fig2}. 
	
\begin{figure}[H]
%\centering
	\includegraphics[width=11.5 cm]{Fig_02.jpg}
\caption{Specimen soil samples. \textbf{Left}: the lower surface of a stabilized specimen. \textbf{Right}: the upper surface of a stabilized~specimen. \label{fig2}}
\end{figure}   

When {compacting} the specimens, separation was observed in raw soil material. During~the fabrication of samples, all the possible precautionary measures were taken in order to minimize separation of soil granula. However, a~certain separation occurred with the effect varied depend on different binders. The~energy and bio ashes demonstrated more separation compared to the others. The~separation of the specimen can be seen in Figure~\ref{fig2} which illustrates the bottom surface of a packed specimen on the left. It indicates an underfilled material with few fine material between the coarser particles. The~right part of Figure~\ref{fig2} presents the upper surface of the wrapped specimen. This figure illustrated the overfilled soil specimen, which means that it includes the excessive amount of the fine material. This results in that fact that coarser particles have no direct contact with each other which significantly affects the cohesion properties of~soil.  

\subsection{Soil Stabilization by~Binders}

The practical goal was that stabilized soil specimens reach the compressive strength of 4 MPa, since it is a common and acceptable value in road construction context. {The compressive testing of soil specimens stabilized by various binder combinations was used to evaluate how the clay material will react when it is being compressed due to the external loads during exploitation of the roads. }The choice of binders was limited to five different variants{: }cement, cement and slag,{ }lime and slag binders. {Two} binder types include the traditional binders (cement and lime) while the other three binders are the alternative novel types: slag GGBFS, energy ash and bio~ash. 

Cement was used as a main reference binder, because~it is a well documented binder widely used for {stabilization of} soil{ used as a material for construction of road pavement}. The~ashes were regarded as the binders from the category of residual products. The~total amount of binder was chosen in such a way that they would be representative for road stabilization. The~selected binder recipes used for soil stabilization were based on the previous experience {to ensure the quality of soil and the lifetime of roads through improved its bearing capacity. Thus, the~improved quality of soil used for road construction improves the lifetime of a road, decreases expensive maintenance activities to repair roads and road condition}. The~significance level of the results was ensured with statistical trial~planning. 

The general aim was to increase the resistance of soil towards the freeze--thaw cycles, provided that there is no continuous water penetration on the road. To~this end, no extra water was added in the freeze--thaw experiments. {The second aim is to evaluate the  behavior of soil and its response under crushing loads, as~well as plastic flow behavior and ductile fracture limits. }The equipment used in this test was originally developed to compress the fine-grained stabilized moraines, and~it worked well for road base applications. For~compression of the stabilized base course gravel, the~equipment works satisfactorily, but~the pore structures affected its performance for the fine-grained soils. During~the compaction, the~separation was reduced to obtain a gap between the two soil layers forming the~specimens.

\subsection{Seismic~Measurements}
After collecting and refining the soil samples, the~specimens were stored for 90 days. During~this {curing} period, the~resonant frequency measurements were performed on the soil specimens using existing methods described earlier as free--free resonate column tests~\citep{Aldaood,Lindh2021b,Lindh2022c,ThomasRangaswamy,Ryden}. The~seismic measurements in this study were performed as evaluation of P-wave velocity after the {curing} of the {soil }specimens on day 14. The~paraffin on the end faces was removed during the actual measurements after which the end faces were resealed. According to this method, the~experiments were repeated with measured both axial values (compression wave) and radial values (shear wave). The~radial measurements gave a large scattering in values caused by the low strength of specimens and performance outside of the plastic~tube.

\subsection{Statistical~Analysis}
After the completion of the previous steps, the~statistical analysis was performed for objective data quality control, as~data might be a subject to errors in course of the experiment. {The scope of the statistical testing was to estimate the best combination of binders in blended mixtures through selecting optimal combination and ratios of binders. The~ultimate goal is to increase the quality of roads and lifetimes of pavements due to the improved soil characteristics. To~this end,} a series of trial tests {was} performed before the final experiment{ }using the statistical hypothesis. The~null hypothesis ($H_0$) in this case is that all the binders give the same strength, while the alternative hypothesis is that they do not give the same strength. Two different assumptions can be made when testing these hypotheses represented by Equations~(\ref{eq01}) and \ref{eq02}:
\begin{equation}\label{eq01}
	\alpha = P (type I wrong) = P(reject H_0 | H_0 is true)
\end{equation}
%\unskip
\vspace{-12pt}
\begin{equation}\label{eq02}
	\beta = P (type II wrong) = P(fails to reject H_0 | H_0 is false)
\end{equation}
where $\alpha$ is usually called the significance level of the test. In~this case, $\alpha = 0.05 $ was used. This means 5\% probability that $H_0$ will be rejected even though the $H_0$ is~true. 

\subsection{Simplex-Lattice~Design}
A mixture design based on the simplex-lattice approach was selected for the experimental setup. A~simplex-lattice design can be adjusted according to either a square model or a special cubic model.  For~a simplex-lattice design with three factors, the~square model is described by the following Equation~(\ref{eq03}):
\begin{equation}\label{eq03}
	y = \beta _1 x_1 +  \beta _2  x_2 +  \beta _3  x_3 + \beta _{1,2} x_1 x_2 + \beta _{1,3} x_1 x_3 + \beta _{2,3} x_2 x_3 + \epsilon
\end{equation}
where $\beta$ is the coefficient of regression; $\epsilon$ is the residual error; $x_1$ is the value of binders which indicate factors; \emph{y} %MDPI: Equation elements should be kept the consistent format with equation. Please check and revise. The same below. <-- Authors' reply: corrected to italics
 is the dependent variable. The~cubic model can be described as follows in the Equation~(\ref{eq04}):
\begin{equation}\label{eq04}
	y = \beta _1 x_1 + \beta _2 x_2 + \beta _3 x_3 + \beta _{1,2} x_1 x_2 + \beta _{1,3} x_1 x_3 + \beta _{2,3} x_2 x_3  + \beta _{1,2,3} x_1 x_2 x_3 + \epsilon
\end{equation}
where $x_1$ is the value of binders, $\beta$ is the coefficient of regression; \emph{y} is the response area and $\epsilon$ describes the error term.  The~difference between the square and cubic model is that in the cubic model, the~cubic term $x_1 x_2 x_3$ is also determined. However, the~increased resolution in the cubic model involves a larger number of experiments which should be performed. In~this study, we defined the following parameters: five different factors corresponding the types of five binders; a cubic design which enables to include an increased number of internal points and an increased proportion of points on the stripes %MDPI: Figures should be numbered in order of appearance. We detected "Figrue 5" appears before "3", please revise and make sure all the figures to appear in numerical order. <-- Authors' reply: yes, agree and confirm. All the figures and their references are checked throughout the text.
and double trials; the total number of specimens is 82 soil~samples.

\subsection{Freeze-Thaw~Tests}
{Fine-grained clayey soils are not suitable as subbase layers in roads and constructions because of their compressibility and frost susceptibility. Therefore, to~improve the engineering properties of fine-grained clayey soils and to increase its shear strength, binders were added and the freeze--thaw experiments were performed to evaluate the effects of stabilization and assess the resistance of soil to reduced temperature below zero. }The freeze--thaw {tests} were performed according to the exiting standard of SS-EN 137244. During~the freeze--thaw experiments{, }soil specimens were frozen for a total of 56 cycles {with no extra water added}. The~{measurements were performed} on days 7, 14, 28, 42 and 56 {with varied} temperature{, as~shown} in Figure~\ref{fig3}. {The freeze--thaw cycles during spring--winter period are typical for northern climate conditions of Sweden. The~freeze--thaw cycles influence the curing of soil and negatively affect the strength of the stabilized soil. Moreover, the~reaction of different soil types varies as a response to freeze--thaw cycles depending to the parameters of stabilization such as curing period, binder content and proportions. Therefore, testing the behavior of stabilized soil during changed temperature is essential for evaluating its future performance prior to road construction}. 

\begin{figure}[H]
%\centering
	\includegraphics[width=8.0 cm]{Fig_03}
\caption{Temperature %MDPI: Please change the hyphen (-) into a minus sign ($-$, "U+2212"). e.g., "-1" should be "$-$1". <-- Authors' reply: Figure 3 is corrected and replotted with hyphen changed to minus sign.
 cycle for the freeze--thaw~experiments. \label{fig3}}
\end{figure}
\unskip 

%%%%%%%%%%%%%%%%%%%%%%%%%%%%%%%%%%%%%%%%%%
\section{Results}

{The impact of the blended binders used for soil stabilization on the engineering properties of clayey soils was investigated with focus on the increase of strength. }We have tested {the effects of }various binders in diverse combination and proportions {on }soil stabilization{. }In this section, we show {the} results {of the improved} strength of the stabilized soil {estimated} by P-wave velocities and visualized {in} ternary diagrams. Upon~{soil stabilisation}, the~obtained values in strength were evaluated statistically and visualized as Pareto charts in Figure~\ref{fig4}. {The results show that the strength increases with curing time and added binders.} 

{Diverse }combination of various binders provides different {effects} on {soil} strength {which prove the} effects from various binders. In~this chart (Figure \ref{fig4}), the~abbreviations for the five binders are denoted by letters: A--cement, B--lime, C--slag GGBFS, D--energy fly ash, E--bio fly ash. Accordingly, the~combinations of binders are given as follows, Figure~\ref{fig4}: BC--%MDPI: We remove the space at both side of ``--'' in this paragraph. Please confirm. <-- Authors' reply: yes, we agree and confirm.
lime+slag, CE--slag+bio fly ash, AC--cement+slag GGBFS, ABC--cement+lime+slag (a triple combination), AE--cement+ bio fly ash, BCE--lime+slag GGBFS BE+bio fly ash, BDE--lime+energy fly ash+bio fly ash, BCD--lime+slag GGBFS+energy fly ash, ACE--cement+slag+bio fly ash, DE--energy fly ash+bio fly ash, and~BD--lime+energy fly~ash.

\begin{figure}[H]
%\centering
\begin{adjustwidth}{-\extralength}{0cm}
\centering
%\centering %% If there is a figure in wide page, please release command \centering
\includegraphics[width=15 cm]{Fig_04}
\end{adjustwidth}
\caption{\textbf{Left}: Pareto diagram of all the effects included in the cubic model for estimated gain in strength of soil on day 14. \textbf{Right}: Pareto diagram of the significant effects from binders on stabilized soil included in the cubic model for the measurements on day 14. %MDPI: Please change the hyphen (-) into a minus sign ($-$, "U+2212"). e.g., "-1" should be "$-$1". <-- Authors' reply: Figure 4 is corrected and replotted with hyphen changed to minus sign.
 \label{fig4}}
\end{figure} 

In general, {the proposed} method tends to find {the }optimal combination of binders with regard to the stabilization of clayey soil where specimens are selected without external and internal undesired pore structure. The~controlling parameters are the strength development of the soil material, final gain in strength and improved resistance of soil to freeze--thaw cycles. Both the destructive testing (uniaxial compressive strength) and non-destructive (seismic) are used as the test methods. Figure~\ref{fig4} presents the significant effects after the 14 days of curing period of the stabilized soil samples. A~visual comparison of the results indicates that the proportions of cement in a blended mixture of binders is still a clearly dominant factor for the compression wave velocity. We next compare the interaction between lime and slag (a mixture of BC, as~explained above) and evaluate their effects on soil strength. Thus, the~combination BC has increased in size and overpassed the bio~ash. 

The P-wave velocity indicating the stabilization effects on soil is demonstrated in ternary diagrams in Figures~\ref{fig5}--\ref{fig7}. In~the evaluation of the experimental results, gradual elimination of non-significant effects has been applied, i.e.,~the first analysis contains all the effects. Then the non-significant effects are removed until the model consists only of significant effects. The~method is called ‘backward elimination’ and described by~\cite{MyersMontgomery}. Figure~\ref{fig5} illustrates the variations of P-wave velocity measured on day 14 during soil stabilization, indicating soil strength in the stabilized specimens by added various types of binders. Further, Figure~\ref{fig5} demonstrates the effects of added cement, energy ash and lime on changed soil strength. Here the left graph displays the increased interaction between lime and slag, while the right graph shows a lesser interaction between cement and energy ash and a larger interaction between cement and slag when stabilising soil~specimens.  

\begin{figure}[H]
%\centering
	\hspace{-8pt}\includegraphics[width=14.0 cm]{Fig_05}
\caption{The P-wave velocity on day 14 during soil stabilization. \textbf{Left}: increased interaction between lime and slag. \textbf{Right}: larger interaction between cement and slag. %MDPI: Please add 0 before decimal dot, e.g., .8728 should be 0.8728. <-- Authors' reply: ok, Figure 5 is corrected and replotted.
 \label{fig5}}
\end{figure} 

To evaluate {the} results {both} qualitatively {and }quantitatively, the~P-wave velocities {were calculated} for soil samples stabilized by various percentage of {the }three binders. Figure~\ref{fig6} (left) compared our results with case of cement, lime and slag binders to those of the energy fly ash and slag binders, Figure~\ref{fig6} (right). The~results containing cement are consistently better with P-wave velocity exceeding 2700 m/s, except~for the blend with high combination of the GGBFS and lime where P-wave velocity exceed 2400 m/s. In~both cases the areas of the highest gain in strength are shown in red colour. The~interaction between cement and energy ash is also illustrated in Figure~\ref{fig6} (left). 

Of the three different residual products,{ }the slag{ }Merox{ }clearly works best. The~ashes work poorly in this stabilization context. Energy ash has a well-documented effect and can harden on its own~\cite{Taylor}. However, this study disclosed that the energy fly ash does not act well as a binder in a stabilization context of clayey soil. One reason for the poor response from the ashes may be{ }low amount of binders which could be below a threshold value. However, a~higher amount of binders would make the use of the product more expensive {which is} currently is not relevant from an economic perspective. In~some cases, other reasons may be the control factors and cause which combination and types of ashes to be selected as components for blended binders. However, the~issue of the permanence should also be taken into account when fabricating blended mixtures. Finally, Figure~\ref{fig6} (right) also indicates that the interaction effects between cement and energy ash are~amplified. 

\begin{figure}[H]
%\centering
	\hspace{-8pt}\includegraphics[width=14.0 cm]{Fig_06}
\caption{The response of the soil strength after compression indicated by the P-wave velocity measured on day 14 upon stabilization. \textbf{Left}: interaction between cement, lime and slag. \textbf{Right}: large interaction between energy fly ash and~slag. %MDPI: Please add 0 before decimal dot, e.g., .8728 should be 0.8728. <-- Authors' reply: ok, Figure 6 is corrected and replotted.
 \label{fig6}}
\end{figure} 

Figure~\ref{fig7} provides another comparison, completing an assessment of the soil strength after stabilization by blended mixture of cement, energy fly ash and bio fly ash binders (left) and the two types of fly ash and lime (right). Thus, Figure~\ref{fig7} (right) reveals the interaction effects between the fly ash and lime. Figure~\ref{fig7} (left) presents the same pattern of the effects from binders on soil stabilization: a negative contribution from the combination ADE (cement/energy ash/bio ash). Both completions indicate the low effects from the energy fly ash on soil stabilization (shown in figures in dark green colour). the strength in using ternary types of diagrams is in its ability to integrate information from triple combination of binders using for soil stabilization, as~illustrated in the presented Figures~\ref{fig5}--\ref{fig7}.

\begin{figure}[H]
%\centering
	\hspace{-8pt}\includegraphics[width=14.0 cm]{Fig_07}
\caption{The response of the soil strength after compression indicated by the P-wave velocity measured on day 14 upon stabilization. \textbf{Left}: increased interactions between cement and fly ash. \textbf{Right}: increased interaction effects between fly ash and~lime. %MDPI: Please add 0 before decimal dot, e.g., .8728 should be 0.8728. <-- Authors' reply: ok, Figure 7 is corrected and replotted.
 \label{fig7}}
\end{figure} 

As the bio fly ash from the SCA Lilla Edet shows poorer results. This may be due to the carbonation of the binder, i.e.,~the presence of lime in ash has reacted with the CO\textsubscript{2} %MDPI: Please confirm if the italics is unnecessary and can be removed. The following highlights are the same. <-- Authors' reply: yes, agree - italics is unnecessary and deleted.
 in the air and inert CaCO\textsubscript{3} (cf. limestone) has been formed. Previous studies arrived with a better effect from the mixture of lime and fly ash in connection with the stabilization of clayey soil~\cite{ANDAVAN20201125,TOKSOZHOZATLIOGLU2021105931,KHADKA2020100327}. There is no clear answer to the question why the energy fly ash only functions as an inert filler. One reason could be that there is a threshold value of the water binder number. If~ash is to be used in larger projects, the~availability of the 'fresh' ash (i.e., recently made product) should be ensured to minimize possible variations in the stabilized road~base.

%%%%%%%%%%%%%%%%%%%%%%%%%%%%%%%%%%%%%%%%%%
\section{Discussion}

{This study verified the feasibility of application of blended binders (cement,  lime, slag, energy fly ash and bio fly ash) on clayey soil stabilization, specifically for road construction in southern Sweden. The~tests included the UCS of soil specimens treated with binder mixtures and measurement of P-wave velocities for evaluation of strength properties. The~results are of importance for road construction and increase the durability and life cycle of roads through the improved construction materials. We} tested five different binders and their effects on stabilization of soil planned for road construction. This is a common factor test {meaning} that {samples } have to {be }evaluated in a {large} strength {diapason}. Furthermore, the~question arises as to how {one} should handle the {amount of }water {in }binder {mixture}. For~the {soil collected in southern Sweden}, the~optimal water ratio is 6\%. With~a varied amount of binder, water binder number will also vary {and affect} the strength of the specimens{. In~contrast, if} water binder number {remains constant}, water ratio {should vary within } the {soil-binder mixture}, which in turn will vary the degree of {compaction} of the specimens. The~degree of {compaction of soil enables} to evaluate the effect of {binders} on {soil} strength.

The {evaluation of soil strength} was performed {using} the applied geophysical methods. Thus, we measured the frequency of the seismic P-waves which notably varied along with changed soil properties: strength, density, stiffness and composition of soil specimens. The~results have been demonstrated on a series of ternary diagrams and two Pareto charts for the comparative analysis. The~quality of the tests was controlled by the standard handbooks and guidances used for the implementation of the soil tests in the Swedish Geotechnical Institute (SGI). {In this study, five different binders were mixed in various proportions with clayey soils with the goal of improving their strength and durability of roads during construction works. Four steps of the laboratory processing of soil samples were performed on unstabilized and stabilized soils, including: (1) soil compaction, (2) soil stabilization by blended binders and UCS test, (3) freeze--thaw durability experiments with varied temperature, and~(4)  analysis of soil strength by seismic measurements using measured velocity of P-waves.}

While {this study uses slag, lime, energy fly ash, bio fly ash and cement} for clay stabilisation, {this} approach can be extended to{ }other alternative binders by adding polymers~\cite{Rezaeimalek} calcium carbide residue (CCR) \cite{Horpibulsuk}, or{ }various lime dosage levels~\cite{Pedarla}. In~addition, gypsum can be used for stabilization of clays to improve the microstructure{ }for the contaminated soils~\cite{Latifi}. Further, the~plasticity properties of the fine-grained clayey soils can be improved by adding traditional binders such as lime and cement{ }added with various content of stabilizers in a mixture~\cite{Mutaz,Rogers} or the combinations of cement and GGBFS~\cite{Andersson}.

{The use of blended binders has been limited in modern soil stabilization practices, although~mixtures of various binders can be regarded as an effective alternative  binders for improving the engineering properties of clayey soil. }In our experiments, the~highest strengths in soil specimens were achieved for the mixtures of slag, lime and cement, while bio fly ash from the SCA Lilla Edet demonstrated poorer results, which may be due to the carbonation of the binder. The~comparable gain in strength is yielded when used in energy fly ash which does not act sufficiently well as a binder in a stabilization context. In~contrast, slag and lime outperformed the soil strength values in terms of the effects from the blended binders on soil. {The results demonstrated that a higher dosage of slag Merox and cement was required to improve the strength of clayey soil. }Thus, a~notable increase in strength of clayey soils was obtained by adding slag, lime and cement with some addition of fly ash. Following the experimental testing, the~study presented a data analysis based on the statistical processing using ternary simplex diagrams for each binder combination and Pareto chart for the comparative~analysis. 

A possible disadvantage of the presented ‘mixture design’ method is that the binder content is not an independent variable, i.e.,~we tested the total amount of binder for each specimen, which is the same but consists of different mixtures: a ratio of single binder components in a binder blend. This means that only a total binder content is tested, compared to the ordinary factorial experiment where two or three levels of binder contents were tested. In~order to treat several different levels of binder content, two or three ‘mixture tests’ must then be performed. As~the study primarily focuses on examining the effects of the akin mixture in an ordinary binder, it was chosen to perform the experiments at the only one level (that is, 3\% mixture). Despite this limitation, 82 different specimens were included in the experiment and tested positively, as~reported~above. 

To address the above drawbacks, we performed the experiments with the constrained design which included the increased amount of the preliminary tests and a more complex evaluation of the results. We applied the third alternative, i.e.,~a mixture experiment, which is a special form of the response surface methodology where factors are the components of the mixture and the response is a function of the proportion of each ingredient. We compared the effects from lime, slag and fly ash binders on clay stabilization. Since this study has only included the inert {soil mass} and storage without the access to air, the~effects of the lime binder were limited. Pure lime (burnt or quenched) is best suited for the fine-grained soils containing clay minerals. Slag is usually referred to as a binder component that requires activation (high pH) to react \citep{Sherwood,Lindh2022a}. 

The specific contribution of this paper consist in using blended binders. Thus, whilst a large number of projects and reported studies utilize pure binders (e.g., cement, lime or~CKD), there is a lack of knowledge regarding the use of mixed combinations of binders for stabilization of weak soils, which experimentally include novel types of binders, such as energy fly ash or bio fly ash. Therefore, we demonstrated that blended binders are particularly effective for stabilising clayey soil in the soft composition. We used these binders in combination with cement, lime and slag for soil stabilization where slag was used with a very finely ground structure and hydraulic properties (hardens under water) without activation. Such properties have also been observed in previous studies \citep{Lindh2000,Lindh2001}. 

{Slag demonstrated to have} a slow curing {effects on clayey soil. Therefore, the~use of pure slag} is not practical{ }in real projects that require {quick stabilization of} dozens of tons of {soil }materials. Therefore, we tested the combination of the blended binders to improve the stabilization of soil. Specifically, the~combination of slag and lime or slag and cement accelerates the process of soil hardening and gives a good and durable product, as~also noted in previous studies~\cite{Lindh2021a}. Hence, this study {contributed to} existing research {gaps} on {evaluation of the effects from }blended binders {on stabilization of }clayey soil {used for }road construction{. Moreover, we demonstrated that using the alternative binders as a partial replacements of cement presents an environmentally friendly method of soil improvement where binders are mixed with clayey soil specimens. Compared to the use of pure cement for soil stabilization, the~use of blended mixtures presents better perspectives for the environmental sustainability and city infrastructure.}

%%%%%%%%%%%%%%%%%%%%%%%%%%%%%%%%%%%%%%%%%%
\section{Conclusions}

We have addressed the problem of using blended binders which are gaining interest in road engineering industry for soil stabilization as efficient materials and environmentally effective replacements to pure cement. This problem is formulated in this paper as optimising the selection of binders using statistical approach and a series of experiments on clay samples. The~research was undertaken to determine and demonstrate how blended binders produce a better effect on clayey soil stabilization, and~to illustrate graphically the performance of soil treated by mixtures of five binders in various combinations: cement, lime, slag, bio fly ash and energy fly ash. The~results were visualized on the ternary simplex diagrams. The~soil strength increased with added blended binders in various ratios, detected using seismic measurements of P-wave~frequency. 

Below we propose a list of further approaches to continue this study to better analyze soil performance stabilized with various~binders:
\begin{itemize}
	\item	Fatigue tests on stabilized material, e.g.,~compare with asphalt. 
	\item Examination of the nonlinearity of soil material (e.g., elongation stiffness)
	\item Testing relationship between seismic, E-modulus and strength of the stabilized soil
	\item Improvement of specimen production
\end{itemize}

Our empirical evaluation of using blended binders demonstrates superior performance of soil and hardening results relative to the state of the art use of single binders. Nevertheless, the~use of blended binders for soil stabilization, specifically novel types, such as energy fly ash or bio fly ash, in~a challenge in the context of road improvement. The~analysis of the experiments determining the effects from various binders on soil behavior performed included simplex lattice design and visualized as ternary plots. The~proposed framework aimed at describing the relations between adding the specific binders on improving soil properties and hardening. Three logical relationships between the ratio of the corresponding five binders were categorized: cement, lime, slag and two types of fly~ash. 

This study revealed that the stabilization of clayey soils with blended binders consisted of lime, fly ash and cement considerably improves the geotechnical properties and workability of soil and contributes to gain in strength. The~significance of this paper is that it also presented a theoretical analysis of the relation between additives and soil properties as estimated by P-wave velocities indicating the stabilization effects on soil. Further, we performed the freeze--thaw experiments for soil collected in Swedish environment and included the statistical analysis aimed at optimization of mixture design. Specifically, we have derived the results from the gain in soil strength after UCS tests for several cases of binders and compared the results in ternary plots. The~results demonstrated the encouraging and positive effects from the use of blended binders, since all the stabilized soil specimens improved the UCS, as~recorded by the P-wave seismic testing. To~conclude, the~bearing capacity of soil was significantly improved by the stabilization using blended binders, which presents the effective solution for the improvement of the engineering materials used for road~constructions.

The key idea of our proposed technique is to use various combinations of binder components in the triplets of mixture blends composed by various stabilizing agents. By~restricting the sum of stabilizing agents always to 100\%, their single components varied in proportions, thus showing the effects of adding more of less of each type of binder on the gain in soil strength. We show that our technique, supported by statistical combinatorics, is able to assist in selecting optimal binder combinations when stabilizing large amount of materials: up to several dozens of tons of soil in the industrial scale during construction works. Our simplex-lattice design is much more efficient than the \emph{in-situ} soil stabilizing due to the computed models showing the effects of  binders on soil strength. It can be easily incorporated into diverse similar construction works where large amount of expansive soil should be stabilized by binder blends. Thus, it can be used as a valuable modelling tool for used for solutions of clay stabilization and evaluation of the degree of gain in soil strength in civil~engineering.
\vspace{6pt}


%%%%%%%%%%%%%%%%%%%%%%%%%%%%%%%%%%%%%%%%%%
\authorcontributions{Supervision, conceptualization, writing---original draft preparation, data curation, methodology, software, visualization, funding acquisition, and~project administration, P.L. (Per Lindh); writing---original draft preparation, methodology, software, formal analysis, resources, validation, writing---review and editing, and~investigation, P.L. (Polina Lemenkova). All authors have read and agreed to the published version of the manuscript.}

\funding{This research was supported by the SBUF (Development Fund of the Swedish Construction Industry), Nordkalk AB, SSAB Merox AB and Svenska Cellulosa AB (SCA) Lilla Edet Pappersbruk. The publication was funded by the %MDPI: Please confirm if the italics is unnecessary and can be removed. <-- Authors' reply: ok, not necessary and removed
Multidisciplinary Digital Publishing Institute (MDPI), by~providing 100\% discount for the APC of this manuscript.}

\institutionalreview{Not applicable.}

\informedconsent{Not applicable.}

\dataavailability{Not applicable.} 

\acknowledgments{The authors thank the anonymous reviewers for reading, suggestions and comments that improved an earlier version of this manuscript.  %MDPI: We suggest  that this sentence should be moved to Funding section. Please revise. <-- Authors' reply: yes, agree and confirm (moved to Funding).
}

\conflictsofinterest{The authors declare no conflict of~interest.} 

\abbreviations{Abbreviations}{
The following abbreviations are used in this manuscript:\\

\noindent 
\begin{tabular}{@{}ll}
GGBFS & Ground Granulated Blast Furnace Slag \\
SGI & Swedish Geotechnical Institute \\
UCS & Uniaxial Compressive Strength 
\end{tabular}
}

%%%%%%%%%%%%%%%%%%%%%%%%%%%%%%%%%%%%%%%%%%
\begin{adjustwidth}{-\extralength}{0cm}
%\printendnotes[custom] % Un-comment to print a list of endnotes

\reftitle{References}

%=====================================
% References, variant A: external bibliography
%=====================================
%\bibliography{ConstrMat.bib}
%\nocite*{}
\begin{thebibliography}{999}

\bibitem[Consoli \em{et~al.}(2019)Consoli, Marin, Samaniego, Heineck, and
  Johann]{Consoli2019}
Consoli, N.C.; Marin, E.J.B.; Samaniego, R.A.Q.; Heineck, K.S.; Johann, A.D.R.
\newblock Use of Sustainable Binders in Soil Stabilization.
\newblock {\em J. Mater. Civ. Eng.} {\bf 2019}, {\em
  31},~06018023.
\newblock {{https://doi.org/10.1061/(ASCE)MT.1943-5533.0002571}}.

\bibitem[da~Rocha \em{et~al.}(2021)da~Rocha, Marin, Samaniego, and
  Consoli]{daRocha}
da~Rocha, C.G.; Marin, E.J.B.; Samaniego, R.A.Q.; Consoli, N.C.
\newblock Decision-Making Model for Soil Stabilization: Minimizing Cost and
  Environmental Impacts.
\newblock {\em J. Mater. Civ. Eng.} {\bf 2021}, {\em
  33},~06020024.
\newblock {{https://doi.org/10.1061/(ASCE)MT.1943-5533.0003551}}.

\bibitem[Freitag \em{et~al.}(2021)Freitag, Maurer, and Reichenauer]{Freitag}
Freitag, P.; Maurer, H.; Reichenauer, T.G.
\newblock Selection of Optimized Binders for Dechlorination of
  Trichloroethylene-Contaminated Sites by Jet Grouting and Deep Soil Mixing.
\newblock {\em J. Environ. Eng.} {\bf 2021}, {\em
  147},~04021052.
\newblock {{https://doi.org/10.1061/(ASCE)EE.1943-7870.0001926}}.

\bibitem[Caraşca(2016)]{Carasca}
Caraşca, O.
\newblock {Soil Improvement by Mixing: Techniques and Performances}.
\newblock {\em Energy Procedia} {\bf 2016}, {\em 85},~85--92.
\newblock https://doi.org/10.1016/\linebreak j.egypro.2015.12.277.

\bibitem[Mypati and Saride(2022)]{Mypati}
Mypati, V.N.K.; Saride, S.
\newblock Feasibility of Alkali-Activated Low-Calcium Fly Ash as a Binder for
  Deep Soil Mixing.
\newblock {\em J. Mater. Civ. Eng.} {\bf 2022}, {\em
  34},~04021410.
\newblock {{https://doi.org/10.1061/(ASCE)MT.1943-5533.0004047}}.

\bibitem[Sargent \em{et~al.}(2012)Sargent, Hughes, Rouainia, and
  Glendinning]{Sargent}
Sargent, P.; Hughes, P.N.; Rouainia, M.; Glendinning, S. Soil Stabilisation
  Using Sustainable Industrial By-Product Binders and Alkali Activation.
\newblock In {\em GeoCongress 2012}; American Society of Civil Engineers: Reston, VA, USA,
  2012; Chapter {GeoCongress 2012: State of the Art and Practice in
  Geotechnical Engineering}; pp. 948--957.
\newblock {{https://doi.org/10.1061/9780784412121.098}}.

\bibitem[Amorim \em{et~al.}(2020)Amorim, Segersson, Körnich, Asker, Olsson,
  and Gidhagen]{Amorim}
Amorim, J.; Segersson, D.; Körnich, H.; Asker, C.; Olsson, E.; Gidhagen, L.
\newblock {High resolution simulation of Stockholm's air temperature and its
  interactions with urban development}.
\newblock {\em Urban Clim.} {\bf 2020}, {\em 32},~100632.
\newblock {{https://doi.org/10.1016/j.uclim.2020.100632}}.

\bibitem[Gustavsson \em{et~al.}(2001)Gustavsson, Bogren, and Green]{Gustavsson}
Gustavsson, T.; Bogren, J.; Green, C.
\newblock {Road climate in cities: A study of the Stockholm area, south-East
  Sweden}.
\newblock {\em Meteorol. Appl.} {\bf 2001}, {\em 8},~481--489.
\newblock {{https://doi.org/10.1017/S1350482701004091}}.

\bibitem[Etim \em{et~al.}(2020)Etim, Attah, and Yohanna]{Etim}
Etim, R.; Attah, I.; Yohanna, P.
\newblock {Experimental study on potential of oyster shell ash in structural
  strength improvement of lateritic soil for road construction}.
\newblock {\em Int. J. Pavement Res. Technol.} {\bf
  2020}, {\em 13},~341--351.
\newblock {{https://doi.org/10.1007/s42947-020-0290-y}}.

\bibitem[Fasihnikoutalab \em{et~al.}(2020)Fasihnikoutalab, Pourakbar, Ball,
  Unluer, and Cristelo]{Fasihnikoutalab}
Fasihnikoutalab, M.H.; Pourakbar, S.; Ball, R.J.; Unluer, C.; Cristelo, N.
\newblock {Sustainable soil stabilization with ground granulated blast-furnace
  slag activated by olivine and sodium hydroxide}.
\newblock {\em Acta Geotech.} {\bf 2020}, {\em 15},~1981--1991.
\newblock {{https://doi.org/10.1007/s11440-019-00884-w}}.

\bibitem[Kogbara \em{et~al.}(2011)Kogbara, Yi, and Al-Tabbaa]{Kogbara}
Kogbara, R.; Yi, Y.; Al-Tabbaa, A.
\newblock {Process envelopes for stabilization/solidification of contaminated
  soil using lime–Slag blend}.
\newblock {\em Environ. Sci. Pollut. Res.} {\bf 2011}, {\em
  18},~1286--1296.
\newblock {{https://doi.org/10.1007/s11356-011-0480-x}}.

\bibitem[Lindh and Lemenkova(2022)]{Lindh2022b}
Lindh, P.; Lemenkova, P.
\newblock {Seismic velocity of P-waves to evaluate strength of stabilized soil
  for Svenska Cellulosa Aktiebolaget Biorefinery Östrand AB, Timrå}.
\newblock {\em Bull. Pol. Acad. Sci. Tech. Sci.}
  {\bf 2022}, {\em 70},~1--9.
\newblock {{https://doi.org/10.24425/bpasts.2022.141593}}.

\bibitem[Mavroulidou \em{et~al.}(2021)Mavroulidou, Gray, Gunn, and
  Pantoja-Muñoz]{Mavroulidou}
Mavroulidou, M.; Gray, C.; Gunn, M.; Pantoja-Muñoz, L.
\newblock {A Study of Innovative Alkali-Activated Binders for Soil
  stabilization in the Context of Engineering Sustainability and Circular
  Economy}.
\newblock {\em Circ. Econ. Sustain.} {\bf 2021}, \emph{2}, 1627--1651.
\newblock {{https://doi.org/10.1007/s43615-021-00112-2}}.

\bibitem[Ekpo \em{et~al.}(2021)Ekpo, Fajobi, and Ayodele]{Ekpo}
Ekpo, D.; Fajobi, A.; Ayodele, A.
\newblock {Response of two lateritic soils to cement kiln dust-periwinkle shell
  ash blends as road sub-base materials}.
\newblock {\em Int. J. Pavement Res. Technol.} {\bf
  2021}, {\em 14},~550--559.
\newblock {{https://doi.org/10.1007/s42947-020-0219-5}}.

\bibitem[Frempong(1995)]{Frempong}
Frempong, E.M.
\newblock {A comparative assessment of sand and lime stabilization of residual
  micaceous compressible soils for road construction}.
\newblock {\em Geotech. Geol. Eng.} {\bf 1995}, {\em
  13},~181--198.
\newblock {{https://doi.org/10.1007/BF00422209}}.

\bibitem[Abbey \em{et~al.}(2021)Abbey, Eyo, Okeke, and Ngambi]{Abbey}
Abbey, S.J.; Eyo, E.U.; Okeke, C.; Ngambi, S.
\newblock {Experimental study on the use of RoadCem blended with by-product
  cementitious materials for stabilization of clay soils}.
\newblock {\em Constr. Build. Mater.} {\bf 2021}, {\em
  280},~122476.
\newblock {{https://doi.org/10.1016/j.conbuildmat.2021.122476}}.

\bibitem[Al-Hdabi \em{et~al.}(2014)Al-Hdabi, {Al Nageim}, and Seton]{AlHdabi}
Al-Hdabi, A.; {Al Nageim}, H.; Seton, L.
\newblock {Superior cold rolled asphalt mixtures using supplementary
  cementations materials}.
\newblock {\em Constr. Build. Mater.} {\bf 2014}, {\em
  64},~95--102.
\newblock {{https://doi.org/10.1016/j.conbuildmat.2014.04.033}}.

\bibitem[Lindh and Lemenkova(2022)]{ma15217798}
Lindh, P.; Lemenkova, P.
\newblock Dynamics of Strength Gain in Sandy Soil Stabilised with Mixed Binders
  Evaluated by Elastic P-Waves during Compressive Loading.
\newblock {\em Materials} {\bf 2022}, {\em 15}, 7798.
\newblock {{https://doi.org/10.3390/ma15217798}}.

\bibitem[Jafer \em{et~al.}(2018)Jafer, Atherton, Sadique, Ruddock, and
  Loffill]{Jafer}
Jafer, H.M.; Atherton, W.; Sadique, M.; Ruddock, F.; Loffill, E.
\newblock {Development of a new ternary blended cementitious binder produced
  from waste materials for use in soft soil stabilization}.
\newblock {\em J. Clean. Prod.} {\bf 2018}, {\em 172},~516--528.
\newblock {{https://doi.org/10.1016/j.jclepro.2017.10.233}}.

\bibitem[Li \em{et~al.}(2021)Li, Chen, Zhan, Wang, Poon, and Tsang]{Lietal2021}
Li, J.S.; Chen, L.; Zhan, B.; Wang, L.; Poon, C.; Tsang, D.
\newblock {Sustainable stabilization/solidification of arsenic-containing soil
  by blast slag and cement blends}.
\newblock {\em Chemosphere} {\bf 2021}, {\em 271},~129868.
\newblock {{https://doi.org/10.1016/j.chemosphere.2021.129868}}.

\bibitem[Wang \em{et~al.}(2015)Wang, Wang, Jin, and Al-Tabbaa]{Wangetal2015}
Wang, F.; Wang, H.; Jin, F.; Al-Tabbaa, A.
\newblock {The performance of blended conventional and novel binders in the
  in-situ stabilization/solidification of a contaminated site soil}.
\newblock {\em J. Hazard. Mater.} {\bf 2015}, {\em 285},~46--52.
\newblock {{https://doi.org/10.1016/j.jhazmat.2014.11.002}}.

\bibitem[Consoli \em{et~al.}(2015)Consoli, Winter, Rilho, Festugato, and {dos
  Santos Teixeira}]{Consoli}
Consoli, N.; Winter, D.; Rilho, A.S.; Festugato, L.; {dos Santos Teixeira}, B.
\newblock {A testing procedure for predicting strength in artificially cemented
  soft soils}.
\newblock {\em Eng. Geol.} {\bf 2015}, {\em 195},~327--334.
\newblock {{https://doi.org/10.1016/j.enggeo.2015.06.005}}.

\bibitem[Meskini \em{et~al.}(2021)Meskini, Samdi, Ejjaouani, and
  Remmal]{MESKINI2021129161}
Meskini, S.; Samdi, A.; Ejjaouani, H.; Remmal, T.
\newblock Valorization of phosphogypsum as a road material: Stabilizing effect
  of fly ash and lime additives on strength and durability.
\newblock {\em J. Clean. Prod.} {\bf 2021}, {\em 323},~129161.
\newblock {{https://doi.org/10.1016/j.jclepro.2021.129161}}.

\bibitem[Mirković \em{et~al.}(2019)Mirković, Tošić, and
  Mladenović]{Mirkovic}
Mirković, K.; Tošić, N.; Mladenović, G.
\newblock {Effect of different types of fly ash on properties of asphalt
  mixtures}.
\newblock {\em Adv. Civ. Eng.} {\bf 2019}, {\em 2019},~1--11.

\bibitem[Navarrete \em{et~al.}(2022)Navarrete, Guimarães, Marques, Castro, and
  Toulkeridis]{app12189134}
Navarrete, C.; Guimarães, A.C.R.; Marques, M.E.S.; Castro, C.D.; Toulkeridis,
  T.
\newblock Resistance to Fatigue in Asphalts Used in Military Airports of the
  Brazilian Amazon through the Use of Nickel-Holding Ash.
\newblock {\em Appl. Sci.} {\bf 2022}, {\em 12}, 9134.
\newblock {{https://doi.org/10.3390/app12189134}}.

\bibitem[{Mayara S. Siverio Lima} \em{et~al.}(2021){Mayara S. Siverio Lima},
  Hajibabaei, Thives, Haritonovs, Buttgereit, Queiroz, and Gschösser]{Lima}
{Lima, M.S.S.}; Hajibabaei, M.; Thives, L.P.; Haritonovs, V.;
  Buttgereit, A.; Queiroz, C.; Gschösser, F.
\newblock Environmental potentials of asphalt mixtures fabricated with red mud
  and fly ash.
\newblock {\em Road Mater. Pavement Des.} {\bf 2021}, {\em
  22},~S690--S701.
\newblock {{https://doi.org/10.1080/14680629.2021.1900899}}.

\bibitem[Bjerrum and Flodin(1960)]{Bjerrum}
Bjerrum, L.; Flodin, N.
\newblock {The Development of Soil Mechanics in Sweden, 1900–1925}.
\newblock {\em Géotechnique} {\bf 1960}, {\em 10},~1--18.
\newblock {{https://doi.org/10.1680/geot.1960.10.1.1}}.

\bibitem[Higgins(2007)]{Higgins}
Higgins, D.
\newblock {Briefing: GGBS and sustainability}.
\newblock {\em Proc. Inst. Civ. Eng.–Constr. Mater.} {\bf 2007}, {\em 160},~99--101.
\newblock https://doi.org/10.1680/\linebreak coma.2007.160.3.99.

\bibitem[Lindh and Lemenkova(2022)]{electronics11223726}
Lindh, P.; Lemenkova, P.
\newblock \textls[-18]{Simplex Lattice Design and X-ray Diffraction for Analysis of Soil
  Structure: A Case of Cement-Stabilised Compacted Tills Reinforced with Steel
  Slag and Slaked Lime.}
\newblock {\em Electronics} {\bf 2022}, {\em 11}, 3726.
\newblock {{https://doi.org/10.3390/electronics11223726}}.

\bibitem[Parsons and Kneebone(2005)]{Parsons}
Parsons, R.; Kneebone, E.
\newblock {Field performance of fly ash stabilized subgrades}.
\newblock {\em Ground Improv.} {\bf 2005}, {\em 9},~33--38.
\newblock {{https://doi.org/10.1680\linebreak /grim.2005.9.1.33}}.

\bibitem[Olgun(2013)]{Olgun}
Olgun, M.
\newblock {Effects of polypropylene fiber inclusion on the strength and volume
  change characteristics of cement-fly ash stabilized clay soil}.
\newblock {\em Geosynth. Int.} {\bf 2013}, {\em 20},~263--275.
\newblock {{https://doi.org/10.1680/gein.13.00016}}.

\bibitem[xin Jie \em{et~al.}(2013)xin Jie, xin Li, Tang, Jin, and xin Hua]{Jie}
Jie, Y.X.; Li, G.X.; Tang, F.; Jin, Y.; Hua, J.X.
\newblock Soil Stabilization in the Fill Project of the Olympic Rowing-Canoeing
  Park in Beijing.
\newblock {\em J. Mater. Civ. Eng.} {\bf 2013}, {\em
  25},~462--471.
\newblock {{https://doi.org/10.1061/(ASCE)MT.1943-5533.0000589}}.

\bibitem[Holz and Batchelor(2012)]{Holz}
Holz, B.; Batchelor, O. Stabilization of Abandoned Coal Mines for a Motorway
  Upgrade, Queensland, Australia.
\newblock In {\em Grouting and Deep Mixing 2012}; American Society of Civil
  Engineers: Reston, VA, USA, 2012; pp. 1912--1923.
\newblock https://doi.org/10.1061/978078441\linebreak2350.0166.

\bibitem[Jamnongpipatkul \em{et~al.}(2009)Jamnongpipatkul, Dechasakulsom, and
  Sukolrat]{Jamnongpipatkul}
Jamnongpipatkul, P.; Dechasakulsom, M.; Sukolrat, J. Application of Air Foam
  Stabilized Soil for Bridge-Embankment Transition Zone in Thailand.
\newblock In {\em Asphalt Material Characterization, Accelerated Testing, and
  Highway Management}; American Society of Civil Engineers: Reston, VA, USA, 2009; pp.
  181--193.
\newblock {{https://doi.org/10.1061/41042(349)24}}.

\bibitem[Moh(1962)]{Moh}
Moh, Z.C.
\newblock Soil Stabilization with Cement and Sodium Additives.
\newblock {\em J. Soil Mech. Found. Div.} {\bf
  1962}, {\em 88},~81--105.
\newblock {{https://doi.org/10.1061/JSFEAQ.0000478}}.

\bibitem[Celauro \em{et~al.}(2012)Celauro, Bevilacqua, Lo~Bosco, and
  Celauro]{Celauro}
Celauro, B.; Bevilacqua, A.; Lo~Bosco, D.; Celauro, C.
\newblock {Design Procedures for Soil-Lime Stabilization for Road and Railway
  Embankments. Part 2-Experimental Validation}.
\newblock {\em Procedia–Soc. Behav. Sci.} {\bf 2012}, {\em
  53},~568--579.
\newblock {{https://doi.org/10.1016/j.sbspro.2012.09.907}}.

\bibitem[Glossop and Wilson(1953)]{Glossop}
Glossop, R.; Wilson, G.C.
\newblock {Soil stability problems in road engineering}.
\newblock {\em Proc. Inst. Civ. Eng.} {\bf 1953},
  {\em 2},~219--253.
\newblock {{https://doi.org/10.1680/ipeds.1953.11550}}.

\bibitem[Markwick and Keep(1942)]{Markwick}
Markwick, A.H.D.; Keep, H.S.
\newblock {The use of low-grade aggregates and soils in the construction of
  bases for roads and aerodromes}.
\newblock {\em Road Eng. Division. Inst. Civ. Eng. Eng. Div. Pap.} {\bf 1942}, {\em 1},~1 –30.
\newblock {{https://doi.org/10.1680/idivp.1942.13348}}.

\bibitem[Nelson and Miller(1992)]{Nelson}
Nelson, J.D.; Miller, D.J.
\newblock {\em Expansive Soils: Problems and Practice in Foundation and
  Pavement Engineering}; John Wiley \& Sons Inc.: Hoboken, NJ, USA, %We added the location of publisher. Please confirm <-- Authors' reply: yes, we agree and confirm.
 1992.

\bibitem[Olsen \em{et~al.}(2022)Olsen, Kasper, and {de. Wit}]{Olsenetal2022}
Olsen, T.; Kasper, T.; {de Wit}, J.
\newblock {Immersed tunnels in soft soil conditions experience from the last 20
  years}.
\newblock {\em Tunn. Undergr. Space Technol.} {\bf 2022}, {\em
  121},~104315.
\newblock {{https://doi.org/10.1016/j.tust.2021.104315}}.

\bibitem[Lindh and Lemenkova(2023)]{LindhLemenkovaTTJ}
Lindh, P.; Lemenkova, P.
\newblock Laboratory Experiments on Soil Stabilization to Enhance Strength
  Parameters for Road Pavement.
\newblock {\em Transp. Telecommun. J.} {\bf 2023}, {\em
  24},~73--82.
\newblock {{https://doi.org/doi:10.2478/ttj-2023-0008}}.

\bibitem[Sear(2001)]{Sear}
Sear, L.K.A. {Properties and use of coal fly ash}.
\newblock In {\em {Use of fly Ash for Road Construction, Runways and Similar
  Projects}}; Thomas Telford Ltd.: London, UK, 2001; chapter~5, pp. 136--168.

\bibitem[Sudhakar \em{et~al.}(2021)Sudhakar, Duraisekaran, {Dilli Vignesh}, and
  {Dhuvarakesh Kanna}]{Sudhakar}
Sudhakar, S.; Duraisekaran, E.; Dilli Vignesh, G.; Kanna, G.D.
\newblock {Performance Evaluation of Quarry Dust Treated Expansive Clay for
  Road Foundations}.
\newblock {\em Iran. J. Sci. Technology. Trans. Civ. Eng.} {\bf 2021}, {\em 45},~2637--2649.
\newblock {{https://doi.org/10.1007/s40996-021-00645-4}}.

\bibitem[Antunes \em{et~al.}(2017)Antunes, Simão, and Freire]{Antunes}
Antunes, V.; Simão, N.; Freire, A.
\newblock {A Soil-Cement Formulation for Road Pavement Base and Sub Base
  Layers: A Case Study}.
\newblock {\em Transp. Infrastruct. Geotechnol.} {\bf 2017}, {\em
  4},~126–141.
\newblock {{https://doi.org/10.1007/s40515-017-0043-9}}.

\bibitem[Park \em{et~al.}(2020)Park, Lee, and Vimonsatit]{Park}
Park, H.K.; Lee, H.; Vimonsatit, V.
\newblock {Investigation of Pindan soil modified with polymer stablisers for
  road pavement}.
\newblock {\em J. Infrastruct. Preserv. Resil.} {\bf
  2020}, {\em 1},~9.
\newblock {{https://doi.org/10.1186/s43065-020-00009-8}}.

\bibitem[Bell(1989)]{Bell}
Bell, F.G.
\newblock {Lime stabilization of clay soils}.
\newblock {\em Bull. Int. Assoc. Eng. Geol.} {\bf 1989}, {\em 39},~67--74.
\newblock {{https://doi.org/10.1007/BF02592537}}.

\bibitem[Hilt and Davidson(1960)]{Hilt}
Hilt, G.; Davidson, D.
\newblock {Lime fixation of clayey soils}.
\newblock {\em Highw. Res. Board Bull.} {\bf 1960}, {\em 262},~20--32.
%\newblock Washington DC, US.

\bibitem[Lindh and Lemenkova(2022)]{LindhLemenkovaCEE}
Lindh, P.; Lemenkova, P.
\newblock Impact of Strength-Enhancing Admixtures on Stabilization of Expansive
  Soil by Addition of Alternative Binders.
\newblock {\em Civ. Environ. Eng.} {\bf 2022}, {\em
  18},~726--735.
\newblock {{https://doi.org/10.2478/cee-2022-0067}}.

\bibitem[Taylor(1997)]{Taylor}
Taylor, H.F.W.
\newblock {\em Cement Chemistry}, 2nd ed.; Thomas Telford Publishing: London, UK,
   1997;
\newblock 361p. {{https://doi.org/10.1680/cc.25929}}.

\bibitem[Okagbue(2007)]{Okagbue}
Okagbue, C.O.
\newblock Stabilization of Clay Using Woodash.
\newblock {\em J. Mater. Civ. Eng.} {\bf 2007}, {\em
  19},~14--18.
\newblock {{https://doi.org/10.1061/(ASCE)0899-1561(2007)19:1(14)}}.

\bibitem[Estabragh \em{et~al.}(2017)Estabragh, Kholoosi, Ghaziani, and
  Javadi]{Estabragh}
Estabragh, A.R.; Kholoosi, M.M.; Ghaziani, F.; Javadi, A.A.
\newblock Stabilization and Solidification of a Clay Soil Contaminated with
  MTBE.
\newblock {\em J. Environ. Eng.} {\bf 2017}, {\em
  143},~04017054.
\newblock {{https://doi.org/10.1061/(ASCE)EE.1943-7870.0001248}}.

\bibitem[Xu and Yi(2021)]{Xu}
Xu, B.; Yi, Y.
\newblock Soft Clay Stabilization Using Three Industry Byproducts.
\newblock {\em J. Mater. Civ. Eng.} {\bf 2021}, {\em
  33},~06021002.
\newblock {{https://doi.org/10.1061/(ASCE)MT.1943-5533.0003710}}.

\bibitem[Sörengård \em{et~al.}(2021)Sörengård, Gago-Ferrero, D.B., and
  Ahrens]{Sorengard}
Sörengård, M.; Gago-Ferrero, P.; Kleja, D.B.; Ahrens, L.
\newblock {Laboratory-scale and pilot-scale stabilization and solidification
  (S/S) remediation of soil contaminated with per- and polyfluoroalkyl
  substances (PFASs)}.
\newblock {\em J. Hazard. Mater.} {\bf 2021}, {\em 402},~123453.
\newblock {{https://doi.org/10.1016/j.jhazmat.2020.123453}}.

\bibitem[Al-Swaidani \em{et~al.}(2016)Al-Swaidani, Hammoud, and
  Meziab]{AlSwaidani}
Al-Swaidani, A.; Hammoud, I.; Meziab, A.
\newblock {Effect of adding natural pozzolana on geotechnical properties of
  lime-stabilized clayey soil}.
\newblock {\em J. Rock Mech. Geotech. Eng.} {\bf
  2016}, {\em 8},~714--725.
\newblock {{https://doi.org/10.1016/j.jrmge.2016.04.002}}.

\bibitem[Feng \em{et~al.}(2021)Feng, Du, Zhou, Zhang, Li, Zhou, and Xia]{Feng}
Feng, Y.S.; Du, Y.J.; Zhou, A.; Zhang, M.; Li, J.S.; Zhou, S.J.; Xia, W.Y.
\newblock {Geoenvironmental properties of industrially contaminated site soil
  solidified/stabilized with a sustainable by-product-based binder}.
\newblock {\em Sci. Total Environ.} {\bf 2021}, {\em 765},~142778.
\newblock {{https://doi.org/10.1016/j.scitotenv.2020.142778}}.

\bibitem[Lindh and Lemenkova(2022)]{LindhLemenkovaECES}
Lindh, P.; Lemenkova, P.
\newblock Leaching of Heavy Metals from Contaminated Soil Stabilised by
  Portland Cement and Slag Bremen.
\newblock {\em Ecol. Chem. Eng. S} {\bf 2022}, {\em
  29},~537--552.
\newblock {{https://doi.org/10.2478/eces-2022-0039}}.

\bibitem[Maheepala \em{et~al.}(2022)Maheepala, Nasvi, Robert, Gunasekara, and
  Kurukulasuriya]{Maheepala}
Maheepala, M.; Nasvi, M.; Robert, D.; Gunasekara, C.; Kurukulasuriya, L.
\newblock {A comprehensive review on geotechnical properties of alkali
  activated binder treated expansive soil}.
\newblock {\em J. Clean. Prod.} {\bf 2022}, {\em 363},~132488.
\newblock {{https://doi.org/10.1016/j.jclepro.2022.132488}}.

\bibitem[Lindh and Lemenkova(2022)]{Lindh2022NL}
Lindh, P.; Lemenkova, P.
\newblock Shear bond and compressive strength of clay stabilised with
  lime/cement jet grouting and deep mixing: A case of Norvik, Nynäshamn.
\newblock {\em Nonlinear Eng.} {\bf 2022}, {\em 11},~693--710.
\newblock {{https://doi.org/10.1515/nleng-2022-0269}}.

\bibitem[Nordmark \em{et~al.}(2022)Nordmark, Vestin, Hansson, and
  Kumpiene]{Nordmark}
Nordmark, D.; Vestin, J.; Hansson, L.; Kumpiene, J.
\newblock {Long-term evaluation of geotechnical and environmental properties of
  ash-stabilized road}.
\newblock {\em J. Environ. Manag.} {\bf 2022}, {\em
  318},~115504.
\newblock {{https://doi.org/10.1016/j.jenvman.2022.115504}}.

\bibitem[Montgomery(1997)]{Montgomery1997}
Montgomery, D.C.
\newblock {\em Design and Analysis of Experiments}; Wiley \& Sons, Inc.: Hoboken, NJ, USA, %We added the location of publisher. Please confirm. <-- Authors' reply: yes, we agree and confirm.
  1997.

\bibitem[Box \em{et~al.}(2005)Box, Hunter, and Hunter]{Box}
Box, G.E.P.; Hunter, W.G.; Hunter, J.S.
\newblock {\em Statistics for Experimenters}, 2nd ed.; Wiley-Interscience: New York, NY, USA, %MDPI: We added the location of publisher. Please confirm. <-- Authors' reply: yes, we agree and confirm.
 2005.
\newblock 655p.

\bibitem[Lindh(2004)]{Lindh2004}
Lindh, P.
\newblock Compaction- and Strength Properties of Stabilized and Unstabilized
  Fine-Grained Tills.
\newblock Ph.D. Thesis, Lund University, Lund, Sweden,  2004.
\newblock {{https://doi.org/10.13140/RG.2.1.1313.6481}}.

\bibitem[Aldaood \em{et~al.}(2021)Aldaood, Khalil, Bouasker, and
  Al-Mukhtar]{Aldaood}
Aldaood, A.; Khalil, A.; Bouasker, M.; Al-Mukhtar, M.
\newblock {Experimental Study on the Mechanical Behavior of Cemented Soil
  Reinforced with Straw Fiber}.
\newblock {\em Geotech. Geol. Eng.} {\bf 2021}, {\em
  39},~2985--3001.
\newblock {{https://doi.org/10.1007/s10706-020-01673-z}}.

\bibitem[Lindh and Lemenkova(2021)]{Lindh2021b}
Lindh, P.; Lemenkova, P.
\newblock {Resonant Frequency Ultrasonic P-Waves for Evaluating Uniaxial
  Compressive Strength of the Stabilized Slag–Cement Sediments}.
\newblock {\em Nord. Concr. Res.} {\bf 2021}, {\em 65},~39--62.
\newblock {{https://doi.org/10.2478/ncr-2021-0012}}.

\bibitem[Lindh and Lemenkova(2022)]{Lindh2022c}
Lindh, P.; Lemenkova, P.
\newblock {Soil contamination from heavy metals and persistent organic
  pollutants (PAH, PCB and HCB) in the coastal area of Västernorrland,
  Sweden}.
\newblock {\em Gospod. Surowcami Miner.–Miner. Resour. Manag.} {\bf 2022}, {\em 38},~147--168.
\newblock {{https://doi.org/10.24425/gsm.2022.141662}}.

\bibitem[Thomas and Rangaswamy(2022)]{ThomasRangaswamy}
Thomas, G.; Rangaswamy, K.
\newblock {Small strain stiffness and site-specific Seismic response of
  nano-silica stabilized soft clay in Kochi, India: A case study}.
\newblock {\em Arab. J. Geosci.} {\bf 2022}, {\em 15},~106.
\newblock {{https://doi.org/10.1007/s12517-021-09315-1}}.

\bibitem[Ryden \em{et~al.}(2006)Ryden, Ekdahl, and Lindh]{Ryden}
Ryden, N.; Ekdahl, U.; Lindh, P.
\newblock {Quality Control of Cement stabilized Soil Using Non-destructive
  Seismic Tests}.
\newblock \emph{Adv. Test. Fresh Cem. Mater. Lect.} \textbf{2006}, \emph{34}, 3--4. %MDPI: \hl{} %MDPI: Newly added information. Please confirm. <-- Authors' reply: yes, we agree and confirm.
%\newblock DGZfP-Proceedings BB 102-CD,  34.

\bibitem[Myers and Montgomery(1995)]{MyersMontgomery}
Myers, R.H.; Montgomery, D.C.
\newblock {\em Response Surface Methodology. Process and Product Optimization
  Using Designed Experiments}; John Wiley \& Sons, Inc.: Hoboken, NJ, USA, %We added the location of publisher. Please confirm. <-- Authors' reply: yes, we agree and confirm.
   1995.

\bibitem[Andavan and Pagadala(2020)]{ANDAVAN20201125}
Andavan, S.; Pagadala, V.K.
\newblock A study on soil stabilization by addition of fly ash and lime.
\newblock {\em Mater. Today: Proc.} {\bf 2020}, {\em 22},~1125--1129.
 %MDPI: Please confirm if it is unnecessary and can be removed. The following highlights are the same. <-- Authors' reply: yes, it can be remoted (we removed it).
 {{https://doi.org/10.1016/j.matpr.2019.11.323}}.

\bibitem[{Toksöz Hozatlıoğlu} and
  Yılmaz(2021)]{TOKSOZHOZATLIOGLU2021105931}
{Toksöz Hozatlıoğlu}, D.; Yılmaz, I.
\newblock Shallow mixing and column performances of lime, fly ash and gypsum on
  the stabilization of swelling soils.
\newblock {\em Eng. Geol.} {\bf 2021}, {\em 280},~105931.
\newblock {{https://doi.org/10.1016/j.enggeo.2020.105931}}.

\bibitem[Khadka \em{et~al.}(2020)Khadka, Jayawickrama, Senadheera, and
  Segvic]{KHADKA2020100327}
Khadka, S.D.; Jayawickrama, P.W.; Senadheera, S.; Segvic, B.
\newblock Stabilization of highly expansive soils containing sulfate using
  metakaolin and fly ash based geopolymer modified with lime and gypsum.
\newblock {\em Transp. Geotech.} {\bf 2020}, {\em 23},~100327.
\newblock {{https://doi.org/10.1016/j.trgeo.2020.100327}}.

\bibitem[Rezaeimalek \em{et~al.}(2018)Rezaeimalek, Nasouri, Huang, and
  Bin-Shafique]{Rezaeimalek}
Rezaeimalek, S.; Nasouri, R.; Huang, J.; Bin-Shafique, S.
\newblock Curing Method and Mix Design Evaluation of a Styrene-Acrylic Based
  Liquid Polymer for Sand and Clay Stabilization.
\newblock {\em J. Mater. Civ. Eng.} {\bf 2018}, {\em
  30},~04018200.
\newblock {{https://doi.org/10.1061/(ASCE)MT.1943-5533.0002396}}.

\bibitem[Horpibulsuk \em{et~al.}(2012)Horpibulsuk, Phetchuay, and
  Chinkulkijniwat]{Horpibulsuk}
Horpibulsuk, S.; Phetchuay, C.; Chinkulkijniwat, A.
\newblock Soil Stabilization by Calcium Carbide Residue and Fly Ash.
\newblock {\em J. Mater. Civ. Eng.} {\bf 2012}, {\em
  24},~184--193.
\newblock {{https://doi.org/10.1061/(ASCE)MT.1943-5533.0000370}}.

\bibitem[Pedarla \em{et~al.}(2010)Pedarla, Chittoori, Puppala, Hoyos, and
  Saride]{Pedarla}
Pedarla, A.; Chittoori, S.; Puppala, A.J.; Hoyos, L.R.; Saride, S. Influence
  of Lime Dosage on Stabilization Effectiveness of Montmorillonite Dominant
  Clays.
\newblock In {\em GeoFlorida 2010}; American Society of Civil Engineers: Reston, VA, USA, 2010;
  pp. 767--776.
\newblock {{https://doi.org/10.1061/41095(365)75}}.

\bibitem[Latifi \em{et~al.}(2018)Latifi, Vahedifard, Siddiqua, and
  Horpibulsuk]{Latifi}
Latifi, N.; Vahedifard, F.; Siddiqua, S.; Horpibulsuk, S.
\newblock Solidification--Stabilization of Heavy Metal--Contaminated Clays
  Using Gypsum: Multiscale Assessment.
\newblock {\em Int. J. Geomech.} {\bf 2018}, {\em
  18},~04018150.
\newblock {{https://doi.org/10.1061/(ASCE)GM.1943-5622.0001283}}.

\bibitem[Mutaz \em{et~al.}(2011)Mutaz, Dafalla, and Shamrani]{Mutaz}
Mutaz, E.; Dafalla, M.A.; Shamrani, M.A., Stabilization of Al-Ghatt Clay Shale
  by Using a Mixture of Lime and Cement.
\newblock In {\em Road Materials and New Innovations in Pavement Engineering};
  American Society of Civil Engineers: Reston, VA, USA, 2011; pp. 81--88.
\newblock {{https://doi.org/10.1061/47634(413)11}}.

\bibitem[Rogers \em{et~al.}(2006)Rogers, Boardman, and Papadimitriou]{Rogers}
Rogers, C.D.; Boardman, D.I.; Papadimitriou, G.
\newblock Stress Path Testing of Realistically Cured Lime and Lime/Cement
  Stabilized Clay.
\newblock {\em J. Mater. Civ. Eng.} {\bf 2006}, {\em
  18},~259--266.
\newblock {{https://doi.org/10.1061/(ASCE)0899-1561(2006)18:2(259)}}.

\bibitem[Andersson \em{et~al.}(2001)Andersson, Carlsson, and
  Leppänen]{Andersson}
Andersson, R.; Carlsson, T.; Leppänen, M. Hydraulic Cement Based Binders for
  Mass Stabilization of Organic Soils.
\newblock In {\em Soft Ground Technology}; American Society of Civil Engineers: Reston, VA, USA, 
   2001; pp. 158--169.
\newblock {{https://doi.org/10.1061/40552(301)13}}.

\bibitem[Sherwood(1993)]{Sherwood}
Sherwood, P.
\newblock {\em Soil Stabilization with Cement and Lime}; TRL, HMSO: Richmond, UK, %MDPI: Newly added information. Please confirm. <-- Authors' reply: yes, agree and confirm.
 1993.

\bibitem[Lindh and Lemenkova(2022)]{Lindh2022a}
Lindh, P.; Lemenkova, P.
\newblock {Geochemical tests to study the effects of cement ratio on potassium
  and TBT leaching and the pH of the marine sediments from the Kattegat Strait,
  Port of Gothenburg, Sweden}.
\newblock {\em Baltica} {\bf 2022}, {\em 35},~47--59.
\newblock {{https://doi.org/10.5200/baltica.2022.1.4}}.

\bibitem[Lindh(2000)]{Lindh2000}
Lindh, P.
\newblock \emph{Soil Stabilization of Fine-Grained Till Soils};
\newblock Technical Report 1008; Lund University: Lund, Sweden, 2000;
\newblock Report TVGT.

\bibitem[Lindh(2001)]{Lindh2001}
Lindh, P.
\newblock {Optimizing binder blends for shallow stabilization of fine-grained
  soils}.
\newblock {\em Ground Improv.} {\bf 2001}, {\em 5},~23--34.
\newblock {{https://doi.org/10.1680/grim.2001.5.1.23}}.

\bibitem[Lindh and Lemenkova(2021)]{Lindh2021a}
Lindh, P.; Lemenkova, P.
\newblock {Evaluation of Different Binder Combinations of Cement, Slag and CKD
  for S/S Treatment of TBT Contaminated Sediments}.
\newblock {\em Acta Mech. Autom.} {\bf 2021}, {\em 15},~236--248.
\newblock {{https://doi.org/10.2478/ama-2021-0030}}.

\end{thebibliography}
\PublishersNote{}

%%%%%%%%%%%%%%%%%%%%%%%%%%%%%%%%%%%%%%%%%%
\end{adjustwidth}
\end{document}
